\begin{thebibliography}{20}
\bibitem{WikiComplexity}
Wikipedia contributors. \emph{Computational complexity theory}.
Subject map requested by the user; consulted 20 November 2026.
\url{https://en.wikipedia.org/wiki/Computational_complexity_theory}.

\bibitem{Cook1971}
Stephen A. Cook. The complexity of theorem-proving procedures.
\emph{Proceedings of the Third Annual ACM Symposium on Theory of Computing}
(1971), 151--158. DOI: \href{https://doi.org/10.1145/800157.805047}{10.1145/800157.805047}.
Author-hosted copy:
\url{https://www.cs.toronto.edu/~sacook/homepage/1971.pdf}.
The iPRA witness and induction audits in this report are our reconstruction.

\bibitem{CookPNP}
Stephen A. Cook. \emph{The P versus NP problem}. Millennium problem
description, 2000.
\url{https://www.cs.toronto.edu/~toni/Courses/Complexity/handouts/cook-clay.pdf}.

\bibitem{Savitch1970}
Walter J. Savitch. Relationships between nondeterministic and deterministic
tape complexities. \emph{Journal of Computer and System Sciences}
4(2) (1970), 177--192.
\href{https://doi.org/10.1016/S0022-0000(70)80006-X}{doi:10.1016/S0022-0000(70)80006-X}.
The publisher abstract was retrieved; the full publisher page denied direct
access. The proof and arithmetic audit here are self-contained.

\bibitem{Immerman1988}
Neil Immerman. Nondeterministic space is closed under complementation.
\emph{SIAM Journal on Computing} 17(5) (1988), 935--938.
Author-hosted copy: \url{https://people.cs.umass.edu/~immerman/pub/space.pdf}.

\bibitem{AroraBarak}
Sanjeev Arora and Boaz Barak. \emph{Computational Complexity: A Modern
Approach}. Cambridge University Press, 2009.
Author's background and draft page:
\url{https://theory.cs.princeton.edu/complexity/}.

\bibitem{Buss1986}
Samuel R. Buss. The polynomial hierarchy and intuitionistic bounded
arithmetic. \emph{Structure in Complexity}, Lecture Notes in Computer
Science 223 (1986), 77--103.
Author's description and article links:
\url{https://mathweb.ucsd.edu/~sbuss/ResearchWeb/intuitionisticBA/}.
This source concerns intuitionistic bounded arithmetic, not an automatic
identification of iPRA-definable functions with polynomial time.

\bibitem{ClayPNP}
Clay Mathematics Institute. \emph{P vs NP}, Millennium Problems.
Listed as unsolved when consulted 20 November 2026.
\url{https://www.claymath.org/millennium/p-vs-np/}.

\bibitem{LeanAxioms}
Lean language reference. \emph{Axioms}; \emph{Validating a Lean Proof}.
\url{https://lean-lang.org/doc/reference/latest/Axioms/}.
The companion project is pinned to Lean 4.19.0; the reference explains
dependency auditing rather than the object-iPRA translation.

\bibitem{ProveItPR}
ProveIt contributors. \emph{ConciseIPRA/Lean/ConciseIPRA/PR.lean},
with \emph{Syntax.lean} and \emph{Deduction.lean}.
Repository commit \texttt{83421593771f877bc1ae98f4229eb35507bff05b},
inspected 20 November 2026.
\url{https://github.com/VladimirReshetnikov/ProveIt/blob/83421593771f877bc1ae98f4229eb35507bff05b/Logic/PrimitiveRecursiveArithmetic/ConciseIPRA/Lean/ConciseIPRA/PR.lean}.

\bibitem{ProveItBounded}
ProveIt contributors. \emph{IPRADerivations/Foundations/BoundedFormulas.lean},
same commit. Object-theory bounded Boolean reflection and witness endpoints.
\url{https://github.com/VladimirReshetnikov/ProveIt/blob/83421593771f877bc1ae98f4229eb35507bff05b/Logic/PrimitiveRecursiveArithmetic/Derivations/Lean/IPRADerivations/Foundations/BoundedFormulas.lean}.

\bibitem{ProveItLeast}
ProveIt contributors. \emph{IPRADerivations/Foundations/LeastBoundedWitness.lean},
same commit.
\url{https://github.com/VladimirReshetnikov/ProveIt/blob/83421593771f877bc1ae98f4229eb35507bff05b/Logic/PrimitiveRecursiveArithmetic/Derivations/Lean/IPRADerivations/Foundations/LeastBoundedWitness.lean}.

\bibitem{ProveItInd}
ProveIt contributors. \emph{IPRADerivations/Foundations/BoundedSyntaxInduction.lean},
same commit. Fixed bounded syntax is lowered to quantifier-free induction.
\url{https://github.com/VladimirReshetnikov/ProveIt/blob/83421593771f877bc1ae98f4229eb35507bff05b/Logic/PrimitiveRecursiveArithmetic/Derivations/Lean/IPRADerivations/Foundations/BoundedSyntaxInduction.lean}.
\end{thebibliography}
